% Source: physical PDF p. 323 (printed p. 300). Coverage: complete Figure
% 29.1 and caption. The contour topology, arrow directions, singularity
% locations, and base point were reconstructed from the full-resolution scan.
\begin{figure}[H]
  \centering
  \begin{tikzpicture}[
    x=1cm,
    y=1cm,
    line cap=round,
    line join=round,
    every node/.style={inner sep=1pt},
    contour/.style={
      draw,
      line width=0.8pt,
      postaction={decorate},
      decoration={markings,
        mark=at position 0.055 with
          {\arrow{Stealth[length=2.0mm,width=1.45mm]}},
        mark=at position 0.285 with
          {\arrow{Stealth[length=2.0mm,width=1.45mm]}},
        mark=at position 0.500 with
          {\arrow{Stealth[length=2.0mm,width=1.45mm]}},
        mark=at position 0.715 with
          {\arrow{Stealth[length=2.0mm,width=1.45mm]}},
        mark=at position 0.945 with
          {\arrow{Stealth[length=2.0mm,width=1.45mm]}}
      }
    }
  ]
    % One continuous counterclockwise contour, deformed inward to the
    % common base point P.
    \draw[contour]
      (0,1.44)
      .. controls (0.20,1.18) and (0.20,-0.91) .. (0.39,-1.14)
      .. controls (0.70,-1.42) and (1.43,-0.78) .. (1.53,0.18)
      .. controls (1.66,1.31) and (0.86,2.08) .. (0,2.12)
      .. controls (-0.87,2.08) and (-1.66,1.30) .. (-1.54,0.14)
      .. controls (-1.43,-0.82) and (-0.71,-1.42) .. (-0.40,-1.14)
      .. controls (-0.21,-0.91) and (-0.21,1.18) .. (0,1.44);

    \fill (0,1.44) circle[radius=2.2pt];
    \node[anchor=south west] at (0.03,1.47) {$P$};

    \fill (-0.86,0.54) circle[radius=1.7pt];
    \fill (0.86,0.54) circle[radius=1.7pt];
    \node[below=2pt] at (-0.86,0.54) {$-u_0$};
    \node[below=2pt] at (0.86,0.54) {$+u_0$};
  \end{tikzpicture}
  \renewcommand{\thefigure}{29.1}
  \caption{Deformation of a contour that circles the complex \(u\) plane
  counterclockwise at large \(\lvert u\rvert\) into a contour that starts at
  a base point \(P\), circles the singularity at \(u=-u_0\)
  counterclockwise back to \(P\), and then circles the singularity at
  \(u=+u_0\) counterclockwise back to \(P\).}
  \label{fig:29.1}
\end{figure}
